% =========================================================
% 3.8.11 Boolean Function Simplification
% Two-column clean layout
% 64 Worked Examples
% =========================================================

\subsection{বুলিয়ান রাশি সরলীকরণ অনুশীলন
(\textenglish{Boolean Function Simplification})}

বুলিয়ান সূত্র ও উপপাদ্য ব্যবহার করে নিচের রাশিগুলো
ধাপে ধাপে সরল করা হয়েছে। প্রতিটি গুরুত্বপূর্ণ ধাপ দেখানো হয়েছে,
যাতে শিক্ষার্থী সমাধানের পদ্ধতি সহজে বুঝতে পারে।

\setlength{\columnsep}{14pt}
\setlength{\columnseprule}{0pt}

\begin{multicols}{2}
\small

% =========================================================
% 1
% =========================================================

\textbf{১.}

\[
\begin{aligned}
F
&=AB+A\overline B+\overline AB\\
&=A(B+\overline B)+\overline AB\\
&=A\cdot1+\overline AB\\
&=A+\overline AB\\
&=A+B
\end{aligned}
\]

{\footnotesize
এখানে \(B+\overline B=1\) এবং
\(A+\overline AB=A+B\) সূত্র ব্যবহার করা হয়েছে।
}

\[
\boxed{F=A+B}
\]

\medskip

% =========================================================
% 2
% =========================================================

\textbf{২.}

\[
\begin{aligned}
F
&=(A+B)(A+C)(\overline A+B)\\
&=(A+B)(\overline A+B)(A+C)\\
&=B(A+C)\\
&=AB+BC
\end{aligned}
\]

{\footnotesize
কারণ,
\[
(X+Y)(\overline X+Y)=Y
\]
}

\[
\boxed{F=AB+BC}
\]

\medskip

% =========================================================
% 3
% =========================================================

\textbf{৩.}

\[
F=(A+B+C)(A+B+\overline C)
\]

ধরি,

\[
X=A+B
\]

তাহলে,

\[
\begin{aligned}
F
&=(X+C)(X+\overline C)\\
&=X+C\overline C\\
&=X+0\\
&=X\\
&=A+B
\end{aligned}
\]

\[
\boxed{F=A+B}
\]

\medskip

% =========================================================
% 4
% =========================================================

\textbf{৪.}

\[
\begin{aligned}
F
&=(A+B)(A+\overline B+C)\\
&=A+B(\overline B+C)\\
&=A+B\overline B+BC\\
&=A+0+BC\\
&=A+BC
\end{aligned}
\]

{\footnotesize
এখানে
\[
(X+Y)(X+Z)=X+YZ
\]
সূত্র ব্যবহার করা হয়েছে।
}

\[
\boxed{F=A+BC}
\]

\medskip

% =========================================================
% 5
% =========================================================

\textbf{৫.}

\[
\begin{aligned}
F
&=AB+A\overline BC+\overline ABC\\
&=A(B+\overline BC)+\overline ABC
\end{aligned}
\]

যেহেতু,

\[
B+\overline BC=B+C
\]

তাই,

\[
\begin{aligned}
F
&=AB+AC+\overline ABC\\
&=AB+AC+BC
\end{aligned}
\]

{\footnotesize
এখানে \(AB+\overline ABC=AB+BC\) ব্যবহার করা হয়েছে।
}

\[
\boxed{F=AB+AC+BC}
\]

\medskip

% =========================================================
% 6
% =========================================================

\textbf{৬.}

\[
\begin{aligned}
F
&=AB+AC+\overline AB\\
&=B(A+\overline A)+AC\\
&=B\cdot1+AC\\
&=B+AC
\end{aligned}
\]

\[
\boxed{F=B+AC}
\]

\medskip

% =========================================================
% 7
% =========================================================

\textbf{৭.}

\[
\begin{aligned}
F
&=AB+A\overline B+AC+\overline AC\\
&=A(B+\overline B)
+C(A+\overline A)\\
&=A\cdot1+C\cdot1\\
&=A+C
\end{aligned}
\]

\[
\boxed{F=A+C}
\]

\medskip

% =========================================================
% 8
% =========================================================

\textbf{৮.}

\[
\begin{aligned}
F
={}&(A+B)(\overline A+C)\\
&\times(B+C)
\end{aligned}
\]

\textenglish{Consensus}-এর দ্বৈত সূত্র অনুযায়ী,

\[
\begin{aligned}
F
&=(A+B)(\overline A+C)\\
&=A\overline A+AC+\overline AB+BC\\
&=AC+\overline AB+BC\\
&=AC+\overline AB
\end{aligned}
\]

\[
\boxed{F=AC+\overline AB}
\]

\medskip

% =========================================================
% 9
% =========================================================

\textbf{৯.}

\[
\begin{aligned}
F
={}&(A+B)(\overline A+B)\\
&\times(A+C)\\
&=B(A+C)\\
&=AB+BC
\end{aligned}
\]

\[
\boxed{F=AB+BC}
\]

\medskip

% =========================================================
% 10
% =========================================================

\textbf{১০.}

\[
\begin{aligned}
F
&=ABC+AB\overline C+A\overline BC\\
&=AB(C+\overline C)+A\overline BC\\
&=AB+A\overline BC\\
&=A(B+\overline BC)\\
&=A(B+C)\\
&=AB+AC
\end{aligned}
\]

\[
\boxed{F=AB+AC}
\]

\medskip

% =========================================================
% 11
% =========================================================

\textbf{১১.}

\[
\begin{aligned}
F
={}&ABC+AB\overline C\\
&+\overline ABC+\overline AB\overline C\\
&=AB(C+\overline C)\\
&\quad+\overline AB(C+\overline C)\\
&=AB+\overline AB\\
&=B(A+\overline A)\\
&=B
\end{aligned}
\]

\[
\boxed{F=B}
\]

\medskip

% =========================================================
% 12
% =========================================================

\textbf{১২.}

\[
\begin{aligned}
F
={}&(A+B+C)(A+B+\overline C)\\
&\times(A+\overline B+C)
\end{aligned}
\]

প্রথম দুইটি গুণনীয়ক থেকে,

\[
(A+B+C)(A+B+\overline C)=A+B
\]

তাই,

\[
\begin{aligned}
F
&=(A+B)(A+\overline B+C)\\
&=A+B(\overline B+C)\\
&=A+BC
\end{aligned}
\]

\[
\boxed{F=A+BC}
\]

\medskip

% =========================================================
% 13
% =========================================================

\textbf{১৩.}

\[
\begin{aligned}
F
&=(A+B)(A+C)(A+D)\\
&=(A+BC)(A+D)\\
&=A+BCD
\end{aligned}
\]

{\footnotesize
এখানে দুইবার
\[
(X+Y)(X+Z)=X+YZ
\]
সূত্র প্রয়োগ করা হয়েছে।
}

\[
\boxed{F=A+BCD}
\]

\medskip

% =========================================================
% 14
% =========================================================

\textbf{১৪.}

\[
\begin{aligned}
F
&=AB+\overline AB+BC+BD\\
&=B(A+\overline A)+BC+BD\\
&=B+BC+BD\\
&=B
\end{aligned}
\]

{\footnotesize
কারণ \(B+BC=B\) এবং \(B+BD=B\)।
}

\[
\boxed{F=B}
\]

\medskip

% =========================================================
% 15
% =========================================================

\textbf{১৫.}

\[
F=AB+\overline AC+BC+AD
\]

এখানে \(BC\) হলো \(AB\) ও \(\overline AC\)-এর
\textenglish{Consensus} পদ।

তাই,

\[
\begin{aligned}
F
&=AB+\overline AC+AD
\end{aligned}
\]

\[
\boxed{F=AB+\overline AC+AD}
\]

\medskip

% =========================================================
% 16
% =========================================================

\textbf{১৬.}

\[
\begin{aligned}
F
={}&(A+B)(\overline A+C)\\
&\times(B+C)(A+D)
\end{aligned}
\]

প্রথম তিনটি গুণনীয়কে
\textenglish{Consensus}-এর দ্বৈত সূত্র প্রয়োগ করে,

\[
\begin{aligned}
F
&=(A+B)(\overline A+C)(A+D)\\
&=(AC+\overline AB)(A+D)\\
&=AC+\overline ABD
\end{aligned}
\]

\[
\boxed{F=AC+\overline ABD}
\]

\medskip

% =========================================================
% 17
% =========================================================

\textbf{১৭.}

\[
\begin{aligned}
F
={}&(A+B)(A+\overline B)\\
&\times(\overline A+C)\\
&=A(\overline A+C)\\
&=A\overline A+AC\\
&=AC
\end{aligned}
\]

\[
\boxed{F=AC}
\]

\medskip

% =========================================================
% 18
% =========================================================

\textbf{১৮.}

\[
\begin{aligned}
F
&=AB+\overline ABC+A\overline C\\
&=B(A+\overline AC)+A\overline C\\
&=B(A+C)+A\overline C\\
&=AB+BC+A\overline C
\end{aligned}
\]

এখানে \(AB\) হলো \(BC\) ও \(A\overline C\)-এর
\textenglish{Consensus} পদ।

তাই,

\[
F=BC+A\overline C
\]

\[
\boxed{F=BC+A\overline C}
\]

\medskip

% =========================================================
% 19
% =========================================================

\textbf{১৯.}

\[
\begin{aligned}
F
={}&(A+B+C)(A+\overline B+C)\\
&\times(\overline A+B+C)
\end{aligned}
\]

প্রথম দুইটি গুণনীয়কের জন্য ধরি,

\[
X=A+C
\]

তাহলে,

\[
(X+B)(X+\overline B)=X
\]

তাই,

\[
\begin{aligned}
F
&=(A+C)(\overline A+B+C)\\
&=C+AB
\end{aligned}
\]

\[
\boxed{F=C+AB}
\]

\medskip

% =========================================================
% 20
% =========================================================

\textbf{২০.}

\[
\begin{aligned}
F
&=A+\overline AB+\overline A\overline BC\\
&=A+B+\overline A\overline BC\\
&=A+B+\overline BC\\
&=A+B+C
\end{aligned}
\]

{\footnotesize
প্রথমে \(A+\overline AB=A+B\),
পরে \(B+\overline BC=B+C\) ব্যবহার করা হয়েছে।
}

\[
\boxed{F=A+B+C}
\]

\medskip

% =========================================================
% 21
% =========================================================

\textbf{২১.}

\[
\begin{aligned}
F
&=(A+B)(\overline A+B)(B+C)\\
&=B(B+C)\\
&=B
\end{aligned}
\]

\[
\boxed{F=B}
\]

\medskip

% =========================================================
% 22
% =========================================================

\textbf{২২.}

\[
\begin{aligned}
F
&=AB+\overline AC+BC+\overline ACD\\
&=AB+\overline AC+BC
\end{aligned}
\]

কারণ,

\[
\overline AC+\overline ACD=\overline AC
\]

এখন \(BC\) হলো \textenglish{Consensus} পদ।

তাই,

\[
F=AB+\overline AC
\]

\[
\boxed{F=AB+\overline AC}
\]

\medskip

% =========================================================
% 23
% =========================================================

\textbf{২৩.}

\[
\begin{aligned}
F
={}&ABC+AB\overline C\\
&+A\overline BC+A\overline B\overline C\\
&=AB(C+\overline C)\\
&\quad+A\overline B(C+\overline C)\\
&=AB+A\overline B\\
&=A(B+\overline B)\\
&=A
\end{aligned}
\]

\[
\boxed{F=A}
\]

\medskip

% =========================================================
% 24
% =========================================================

\textbf{২৪.}

\[
\begin{aligned}
F
={}&(A+B+C)(A+\overline B+C)\\
&\times(\overline A+B+C)\\
&\times(\overline A+\overline B+C)
\end{aligned}
\]

প্রথম দুইটি গুণনীয়ক থেকে,

\[
A+C
\]

এবং শেষ দুইটি থেকে,

\[
\overline A+C
\]

তাই,

\[
\begin{aligned}
F
&=(A+C)(\overline A+C)\\
&=C
\end{aligned}
\]

\[
\boxed{F=C}
\]

\medskip

% =========================================================
% 25
% =========================================================

\textbf{২৫.}

\[
\begin{aligned}
F
&=AB+\overline AB+AC+\overline AD\\
&=B(A+\overline A)+AC+\overline AD\\
&=B+AC+\overline AD
\end{aligned}
\]

\[
\boxed{F=B+AC+\overline AD}
\]

\medskip

% =========================================================
% 26
% =========================================================

\textbf{২৬.}

\[
\begin{aligned}
F
={}&(A+B)(A+\overline B+C)\\
&\times(\overline A+C)\\
&=(A+BC)(\overline A+C)\\
&=AC+\overline ABC+BC\\
&=AC+BC\\
&=C(A+B)
\end{aligned}
\]

\[
\boxed{F=C(A+B)}
\]

\medskip

% =========================================================
% 27
% =========================================================

\textbf{২৭.}

\[
\begin{aligned}
F
={}&AB+\overline ABC\\
&+A\overline BC+\overline A\overline BC
\end{aligned}
\]

\(C\)-যুক্ত পদগুলো একত্র করলে,

\[
\begin{aligned}
F
&=AB+C(\overline AB+A\overline B
+\overline A\overline B)\\
&=AB+C(\overline A+\overline B)
\end{aligned}
\]

এখন,

\[
AB+\overline A+\overline B=1
\]

তাই,

\[
F=AB+C
\]

\[
\boxed{F=AB+C}
\]

\medskip

% =========================================================
% 28
% =========================================================

\textbf{২৮.}

\[
\begin{aligned}
F
={}&(A+B+C)(A+B+\overline C)\\
&\times(A+\overline B+C)\\
&\times(A+\overline B+\overline C)
\end{aligned}
\]

প্রথম দুইটি গুণনীয়ক থেকে,

\[
A+B
\]

শেষ দুইটি থেকে,

\[
A+\overline B
\]

তাই,

\[
\begin{aligned}
F
&=(A+B)(A+\overline B)\\
&=A
\end{aligned}
\]

\[
\boxed{F=A}
\]

\medskip

% =========================================================
% 29
% =========================================================

\textbf{২৯.}

\[
F=AB+\overline AC+BC+\overline AD
\]

এখানে \(BC\) হলো \(AB\) ও \(\overline AC\)-এর
\textenglish{Consensus} পদ।

তাই,

\[
F=AB+\overline AC+\overline AD
\]

\[
\boxed{F=AB+\overline AC+\overline AD}
\]

\medskip

% =========================================================
% 30
% =========================================================

\textbf{৩০.}

\[
\begin{aligned}
F
={}&(A+B)(\overline A+C)(B+C)\\
&\times(A+D)(C+D)
\end{aligned}
\]

প্রথম তিনটি গুণনীয়কে
\textenglish{Consensus}-এর দ্বৈত সূত্র প্রয়োগ করে,

\[
\begin{aligned}
F
&=(A+B)(\overline A+C)(A+D)(C+D)\\
&=(AC+\overline AB)(A+D)(C+D)\\
&=AC+\overline ABD
\end{aligned}
\]

\[
\boxed{F=AC+\overline ABD}
\]

\medskip

% =========================================================
% 31
% =========================================================

\textbf{৩১.}

\[
\begin{aligned}
F
={}&AB+\overline ABC+A\overline BC\\
&+\overline A\overline BC+AD
\end{aligned}
\]

প্রথম চারটি পদ সরল করলে,

\[
AB+\overline ABC
+A\overline BC+\overline A\overline BC
=
AB+C
\]

তাই,

\[
F=AB+C+AD
\]

\[
\boxed{F=C+AB+AD}
\]

\medskip

% =========================================================
% 32
% =========================================================

\textbf{৩২.}

\[
\begin{aligned}
F
={}&ABCD+ABC\overline D\\
&+A\overline BCD+A\overline BC\overline D\\
&=ABC(D+\overline D)\\
&\quad+A\overline BC(D+\overline D)\\
&=ABC+A\overline BC\\
&=AC(B+\overline B)\\
&=AC
\end{aligned}
\]

\[
\boxed{F=AC}
\]

\medskip

% =========================================================
% 33
% =========================================================

\textbf{৩৩.}

\[
\begin{aligned}
F
={}&ABCD+AB\overline CD\\
&+A\overline BCD+\overline ABCD\\
&=D(ABC+AB\overline C\\
&\qquad+A\overline BC+\overline ABC)
\end{aligned}
\]

ভেতরের রাশি সরল করলে,

\[
ABC+AB\overline C
+A\overline BC+\overline ABC
=
AB+AC+BC
\]

তাই,

\[
\boxed{F=ABD+ACD+BCD}
\]

\medskip

% =========================================================
% 34
% =========================================================

\textbf{৩৪.}

\[
\begin{aligned}
F
={}&(A+B+C+D)\\
&\times(A+B+C+\overline D)\\
&\times(A+B+\overline C+D)\\
&\times(A+\overline B+C+D)
\end{aligned}
\]

প্রথম দুইটি গুণনীয়ক থেকে,

\[
A+B+C
\]

এরপর ধাপে ধাপে সরল করলে,

\[
F=A+BC+BD+CD
\]

\[
\boxed{F=A+BC+BD+CD}
\]

\medskip

% =========================================================
% 35
% =========================================================

\textbf{৩৫.}

\[
F=AB+\overline AC+BC+\overline AD+CD
\]

এখানে \(BC\) হলো \(AB\) ও \(\overline AC\)-এর
\textenglish{Consensus} পদ।

তাই,

\[
F=AB+\overline AC+\overline AD+CD
\]

\[
\boxed{F=AB+\overline AC+\overline AD+CD}
\]

\medskip

% =========================================================
% 36
% =========================================================

\textbf{৩৬.}

\[
\begin{aligned}
F
={}&(A+B)(\overline A+C)(B+C)\\
&\times(\overline B+D)(C+D)
\end{aligned}
\]

প্রথম তিনটি গুণনীয়ক থেকে,

\[
(A+B)(\overline A+C)
\]

তাই,

\[
\begin{aligned}
F
={}&(AC+\overline AB)\\
&\times(\overline B+D)(C+D)
\end{aligned}
\]

বণ্টন ও শোষণ সূত্র প্রয়োগ করলে,

\[
\boxed{F=BCD+A\overline BC+\overline ABD}
\]

\medskip

% =========================================================
% 37
% =========================================================

\textbf{৩৭.}

\[
\begin{aligned}
F
={}&ABC+AB\overline C\\
&+\overline ABC+\overline A\overline BD\\
&=AB(C+\overline C)\\
&\quad+\overline ABC+\overline A\overline BD\\
&=AB+\overline ABC+\overline A\overline BD\\
&=AB+BC+\overline A\overline BD
\end{aligned}
\]

\[
\boxed{F=AB+BC+\overline A\overline BD}
\]

\medskip

% =========================================================
% 38
% =========================================================

\textbf{৩৮.}

\[
\begin{aligned}
F
={}&(A+B+C)(A+\overline B+C)\\
&\times(\overline A+B+D)\\
&\times(\overline A+\overline B+D)
\end{aligned}
\]

প্রথম দুইটি গুণনীয়ক থেকে,

\[
A+C
\]

শেষ দুইটি থেকে,

\[
\overline A+D
\]

তাই,

\[
\begin{aligned}
F
&=(A+C)(\overline A+D)\\
&=AD+\overline AC+CD\\
&=AD+\overline AC
\end{aligned}
\]

এখানে \(CD\) হলো \textenglish{Consensus} পদ।

\[
\boxed{F=AD+\overline AC}
\]

\medskip

% =========================================================
% 39
% =========================================================

\textbf{৩৯.}

\[
\begin{aligned}
F
={}&AB+A\overline BC\\
&+\overline ABC+\overline A\overline BD
\end{aligned}
\]

প্রথম তিনটি পদ সরল করলে,

\[
AB+A\overline BC+\overline ABC
=
AB+AC+BC
\]

তাই,

\[
\boxed{F=AB+AC+BC+\overline A\overline BD}
\]

\medskip

% =========================================================
% 40
% =========================================================

\textbf{৪০.}

\[
\begin{aligned}
F
&=\overline{A+B}+AB\\
&=\overline A\,\overline B+AB
\end{aligned}
\]

দুইটি পদ পরস্পরকে শোষণ করে না।

তাই এটিই সরল রূপ।

\[
\boxed{F=\overline A\,\overline B+AB}
\]

\medskip

% =========================================================
% 41
% =========================================================

\textbf{৪১.}

\[
\begin{aligned}
F
&=\overline{AB}+\overline AB\\
&=(\overline A+\overline B)+\overline AB\\
&=\overline A+\overline B
\end{aligned}
\]

{\footnotesize
কারণ,
\[
\overline A+\overline AB=\overline A
\]
}

\[
\boxed{F=\overline A+\overline B}
\]

\medskip

% =========================================================
% 42
% =========================================================

\textbf{৪২.}

\[
\begin{aligned}
F
&=\overline{A+B}+\overline AB+A\overline B\\
&=\overline A\,\overline B+\overline AB
+A\overline B\\
&=\overline A(\overline B+B)+A\overline B\\
&=\overline A+A\overline B\\
&=\overline A+\overline B
\end{aligned}
\]

\[
\boxed{F=\overline A+\overline B}
\]

\medskip

% =========================================================
% 43
% =========================================================

\textbf{৪৩.}

\[
F=\overline{AB}+AB
\]

ধরি,

\[
X=AB
\]

তাহলে,

\[
\begin{aligned}
F
&=\overline X+X\\
&=1
\end{aligned}
\]

\[
\boxed{F=1}
\]

\medskip

% =========================================================
% 44
% =========================================================

\textbf{৪৪.}

\[
\begin{aligned}
F
&=\overline{A+B+C}
+\overline A\,\overline BC\\
&=\overline A\,\overline B\,\overline C
+\overline A\,\overline BC\\
&=\overline A\,\overline B(\overline C+C)\\
&=\overline A\,\overline B
\end{aligned}
\]

\[
\boxed{F=\overline A\,\overline B}
\]

\medskip

% =========================================================
% 45
% =========================================================

\textbf{৪৫.}

\[
\begin{aligned}
F
&=\overline{A+\overline{(B+C)}}\\
&=\overline A\,
\overline{\overline{(B+C)}}\\
&=\overline A(B+C)
\end{aligned}
\]

{\footnotesize
এখানে \textenglish{De Morgan's Theorem}
এবং দ্বৈত পূরক সূত্র ব্যবহার করা হয়েছে।
}

\[
\boxed{F=\overline A(B+C)}
\]

\medskip

% =========================================================
% 46
% =========================================================

\textbf{৪৬.}

\[
\begin{aligned}
F
&=\overline{\overline{(A+B)}+C}\\
&=\overline{\overline{(A+B)}}\,\overline C\\
&=(A+B)\overline C
\end{aligned}
\]

{\footnotesize
এখানে \textenglish{De Morgan's Theorem}
এবং \(\overline{\overline X}=X\) ব্যবহার করা হয়েছে।
}

\[
\boxed{F=(A+B)\overline C}
\]

\medskip

% =========================================================
% 47
% =========================================================

\textbf{৪৭.}

\[
\begin{aligned}
F
&=\overline{(A+B)(A+C)}\\
&=\overline{A+B}+\overline{A+C}\\
&=\overline A\,\overline B
+\overline A\,\overline C\\
&=\overline A(\overline B+\overline C)
\end{aligned}
\]

\[
\boxed{F=\overline A(\overline B+\overline C)}
\]

\medskip

% =========================================================
% 48
% =========================================================

\textbf{৪৮.}

\[
\begin{aligned}
F
&=\overline{AB+\overline AC}\\
&=\overline{AB}\,
\overline{\overline AC}\\
&=(\overline A+\overline B)(A+\overline C)
\end{aligned}
\]

বণ্টন করলে,

\[
\begin{aligned}
F
&=\overline A\,\overline C
+A\overline B
+\overline B\,\overline C
\end{aligned}
\]

এখানে \(\overline B\,\overline C\) হলো
\textenglish{Consensus} পদ।

তাই,

\[
\boxed{F=\overline A\,\overline C+A\overline B}
\]

\medskip

% =========================================================
% 49
% =========================================================

\textbf{৪৯.}

\[
\begin{aligned}
F
&=\overline{(A+B)(\overline A+C)}\\
&=\overline{A+B}
+\overline{\overline A+C}\\
&=\overline A\,\overline B+A\overline C
\end{aligned}
\]

\[
\boxed{F=\overline A\,\overline B+A\overline C}
\]

\medskip

% =========================================================
% 50
% =========================================================

\textbf{৫০.}

\[
\begin{aligned}
F
&=\overline{A+B\overline C}\\
&=\overline A\,\overline{B\overline C}\\
&=\overline A(\overline B+C)
\end{aligned}
\]

\[
\boxed{F=\overline A(\overline B+C)}
\]

\medskip

% =========================================================
% 51
% =========================================================

\textbf{৫১.}

\[
F=\overline{AB+\overline AC+BC}
\]

ভেতরের \(BC\) হলো \textenglish{Consensus} পদ।

তাই,

\[
\begin{aligned}
F
&=\overline{AB+\overline AC}\\
&=(\overline A+\overline B)(A+\overline C)\\
&=\overline A\,\overline C+A\overline B
\end{aligned}
\]

\[
\boxed{F=\overline A\,\overline C+A\overline B}
\]

\medskip

% =========================================================
% 52
% =========================================================

\textbf{৫২.}

\[
F=
\overline{(A+B+C)(A+B+\overline C)}
\]

কারণ,

\[
(A+B+C)(A+B+\overline C)=A+B
\]

তাই,

\[
\begin{aligned}
F
&=\overline{A+B}\\
&=\overline A\,\overline B
\end{aligned}
\]

\[
\boxed{F=\overline A\,\overline B}
\]

\medskip

% =========================================================
% 53
% =========================================================

\textbf{৫৩.}

\[
F=
\overline{(A+B)(A+\overline B+C)}
\]

ভেতরের রাশি,

\[
(A+B)(A+\overline B+C)=A+BC
\]

তাই,

\[
\begin{aligned}
F
&=\overline{A+BC}\\
&=\overline A\,\overline{BC}\\
&=\overline A(\overline B+\overline C)
\end{aligned}
\]

\[
\boxed{F=\overline A(\overline B+\overline C)}
\]

\medskip

% =========================================================
% 54
% =========================================================

\textbf{৫৪.}

\[
\begin{aligned}
F
&=\overline{AB+\overline AB}\\
&=\overline{B(A+\overline A)}\\
&=\overline{B\cdot1}\\
&=\overline B
\end{aligned}
\]

\[
\boxed{F=\overline B}
\]

\medskip

% =========================================================
% 55
% =========================================================

\textbf{৫৫.}

\[
\begin{aligned}
F
={}&\overline{
ABC+A\overline BC
+\overline ABC
+\overline A\overline BC
}
\end{aligned}
\]

ভেতরে \(C\) সাধারণ গুণনীয়ক নিলে,

\[
\begin{aligned}
F
&=\overline{
C(AB+A\overline B
+\overline AB+\overline A\overline B)
}\\
&=\overline{C\cdot1}\\
&=\overline C
\end{aligned}
\]

\[
\boxed{F=\overline C}
\]

\medskip

% =========================================================
% 56
% =========================================================

\textbf{৫৬.}

\[
F=\overline{(A+B)(\overline A+B)}
\]

কারণ,

\[
(A+B)(\overline A+B)=B
\]

তাই,

\[
F=\overline B
\]

\[
\boxed{F=\overline B}
\]

\medskip

% =========================================================
% 57
% =========================================================

\textbf{৫৭.}

ধরি,
\[\begin{aligned}
P={}&(A+B+C)(A+B+\overline C)\\
&\times(A+\overline B+C)(A+\overline B+\overline C)
\end{aligned}\]
তাহলে প্রদত্ত রাশি \(F=\overline P\)।

ভেতরের প্রথম দুইটি গুণনীয়ক থেকে,

\[
A+B
\]

এবং শেষ দুইটি থেকে,

\[
A+\overline B
\]

তাই,

\[
\begin{aligned}
F
&=\overline{(A+B)(A+\overline B)}\\
&=\overline A
\end{aligned}
\]

\[
\boxed{F=\overline A}
\]

\medskip

% =========================================================
% 58
% =========================================================

\textbf{৫৮.}

\[
\begin{aligned}
F
&=\overline{A\overline B+\overline AB}\\
&=\overline{A\overline B}\,
\overline{\overline AB}\\
&=(\overline A+B)(A+\overline B)
\end{aligned}
\]

বণ্টন করলে,

\[
\begin{aligned}
F
&=AB+\overline A\,\overline B
\end{aligned}
\]

\[
\boxed{F=AB+\overline A\,\overline B}
\]

\medskip

% =========================================================
% 59
% =========================================================

\textbf{৫৯.}

\[
F=
\overline{(A+B)(\overline A+C)(B+C)}
\]

\textenglish{Consensus}-এর দ্বৈত সূত্র অনুযায়ী,

\[
(A+B)(\overline A+C)(B+C)
=
(A+B)(\overline A+C)
\]

তাই,

\[
\begin{aligned}
F
&=\overline{(A+B)(\overline A+C)}\\
&=\overline{AC+\overline AB}\\
&=(\overline A+\overline C)(A+\overline B)\\
&=\overline A\,\overline B+A\overline C
\end{aligned}
\]

\[
\boxed{F=\overline A\,\overline B+A\overline C}
\]

\medskip

% =========================================================
% 60
% =========================================================

\textbf{৬০.}

\[
\begin{aligned}
F
&=\overline{A+\overline AB}\\
&=\overline{A+B}\\
&=\overline A\,\overline B
\end{aligned}
\]

{\footnotesize
প্রথমে \(A+\overline AB=A+B\),
এরপর \textenglish{De Morgan's Theorem} প্রয়োগ করা হয়েছে।
}

\[
\boxed{F=\overline A\,\overline B}
\]

\medskip

% =========================================================
% 61
% =========================================================

\textbf{৬১.}

\[
\begin{aligned}
F
&=\overline{A(\overline A+B)}\\
&=\overline{A\overline A+AB}\\
&=\overline{0+AB}\\
&=\overline{AB}\\
&=\overline A+\overline B
\end{aligned}
\]

\[
\boxed{F=\overline A+\overline B}
\]

\medskip

% =========================================================
% 62
% =========================================================

\textbf{৬২.}

\[
F=\overline{(A+B)(A+C)(A+D)}
\]

ভেতরের রাশি,

\[
(A+B)(A+C)(A+D)=A+BCD
\]

তাই,

\[
\begin{aligned}
F
&=\overline{A+BCD}\\
&=\overline A\,\overline{BCD}\\
&=\overline A
(\overline B+\overline C+\overline D)
\end{aligned}
\]

\[
\boxed{
F=\overline A
(\overline B+\overline C+\overline D)
}
\]

\medskip

% =========================================================
% 63
% =========================================================

\textbf{৬৩.}

\[
F=
\overline{AB+\overline AC+BC+AD}
\]

ভেতরের \(BC\) হলো \(AB\) ও \(\overline AC\)-এর
\textenglish{Consensus} পদ।

তাই,

\[
F=\overline{AB+\overline AC+AD}
\]

এখন \textenglish{De Morgan's Theorem} প্রয়োগ করে,

\[
\begin{aligned}
F
={}&(\overline A+\overline B)\\
&\times(A+\overline C)
(\overline A+\overline D)
\end{aligned}
\]

অতএব,

\[
\boxed{
\begin{aligned}
F
={}&(\overline A+\overline B)
(A+\overline C)\\
&\times(\overline A+\overline D)
\end{aligned}
}
\]

\medskip

% =========================================================
% 64
% =========================================================

\textbf{৬৪.}

ধরি,
\[\begin{aligned}
P={}&(A+B+C)\\
&\times(A+B+\overline C)(A+\overline B+C)
\end{aligned}\]
তাহলে প্রদত্ত রাশি \(F=\overline P\)।

প্রথম দুইটি গুণনীয়ক থেকে,

\[
(A+B+C)(A+B+\overline C)=A+B
\]

তাই ভেতরের রাশি,

\[
\begin{aligned}
&(A+B)(A+\overline B+C)\\
&=A+B(\overline B+C)\\
&=A+BC
\end{aligned}
\]

অতএব,

\[
\begin{aligned}
F
&=\overline{A+BC}\\
&=\overline A\,\overline{BC}\\
&=\overline A(\overline B+\overline C)
\end{aligned}
\]

\[
\boxed{
F=\overline A(\overline B+\overline C)
}
\]

\end{multicols}

% Restore column rule after this section if needed
\setlength{\columnseprule}{0pt}